\begin{frame}{자주 묻는 질문 (5.2) --- 실제 블랙잭에 쓸 수
있나, Q 부호 해석}
  \begin{itemize}
    \item \textbf{학습한 정책을 실제 블랙잭에 쓰면 이길 수
      있나요?} --- 이 절의 단순화 모델(소프트 에이스 무시,
      무한 카드 덱)로 학습한 정책은 \textbf{기본 전략의
      골격}은 재현하지만, 실제 카지노 규칙(소프트 17,
      멀티덱, 히트 온 17 등)에는 정확히 맞지 않는다. MC 의
      강점은 ``규칙이 아무리 복잡해도(전이확률이 계산
      불가능해도) 시뮬레이션만 하면 전략을 배운다''는 것이지,
      ``학습한 정책이 무조건 현역 프로보다 낫다''는 것이
      아니다.
    \item \textbf{$Q$ 값의 부호(+/−)는 어떻게 해석하나?} ---
      블랙잭은 (공정한 게임이므로) 전체 기대 수익이 0 근처
      이므로, 모든 상태·행동의 $Q$ 가 동시에 양수/음수일 수
      없고 \textbf{상태마다 양/음이 섞여 있어야 한다}. 학습된
      $Q$ 에서 (20, 3) 스틱 $+0.63$, (12, 10) 스틱 $-0.54$ 로
      부호가 반대인 것이 바로 이 ``게임 전체는 0 근처''
      구조를 반영한다.
  \end{itemize}
\end{frame}
